\section[Beilinson--Drinfeld Grassmann 簇]{Beilinson--Drinfeld Grassmann 簇（1 月 29 日）}

\srcpage{124}

\subsection{Beilinson--Drinfeld Grassmann 簇}

假设 $G/\cO_E$ 是分裂约化群。（一般情形下，用 étale 局部化化归到这一情形。\footnote{这是因为 Grassmann 簇本质上是局部对象。对于整体对象，例如 Fargues--Fontaine 曲线，就不能这样做。}）

令 $S=\Spa(R,R^+)$，$\varpi\in R$ 是伪一致化子，$\cY_S=\Spa W_{\cO_E}(R^+)\setminus\{[\varpi]=0\}$。\footnote{先前我们主要在 $\cY_S$ 上工作，那里还设 $\pi\neq0$；但本部分我们要保留向特征 $p$ 的退化。}

回忆我们在 $\cY_S=S\times\Spa\cO_E$ 上有 $d$ 次 Cartier 除子的模空间
\[
\Div^d_{\cY}=(\Div^1_{\cY})^d/\Sigma_d=(\Spa\cO_E)^{\diamond,d}/\Sigma_d.
\]
这是一个小 $v$-叠，且 $\Div^d_{\cY}\to*$ 在局部空间菱形中可表。它「参数化 $\Spa\cO_E$ 上的 $d$ 个点」。

给定 $S$ 以及 $\Div^d_{\cY}(S)$ 的一个截面，我们得到相对 Cartier 除子 $D_S\subset\cY_S$，其理想层为 $\cI(D_S)$。若 $S=\Spa(R,R^+)$ 是 affinoid 的，令 $B^+$ 为 $\cO(\cY_S)$ 沿 $\cI(D_S)$ 的完备化。换言之，
\[
B^+:=W_{\cO_E}(R^+)[1/[\varpi]]^\wedge_\xi,
\]
其中 $D_S=V(\xi)$。定义 $B:=B^+[1/\xi]$。（在 $d=1$ 的情形，它们先前叫作 $B^+_{\mathrm{dR}}$ 与 $B_{\mathrm{dR}}$。）粗略地说，「$B=\cO((\cY_S)^\wedge_{D_S}\setminus D_S)$」。

\begin{warning}
这与 [FF] 中出现的记号 $B,B^+$ 不同。
\end{warning}

\begin{definition}
我们定义 $L^+G/\Div^d_{\cY}$ 为（$\Div^d_{\cY}$ 上的）小 $v$-层
\[
S=\Spa(R,R^+)/\Div^d_{\cY}\mapsto G(B^+),
\]
并定义 $LG/\Div^d_{\cY}$ 为（$\Div^d_{\cY}$ 上的）小 $v$-层
\[
S=\Spa(R,R^+)/\Div^d_{\cY}\mapsto G(B).
\]
我们定义 \emph{Beilinson--Drinfeld Grassmann 簇}
\[
\Gr_{G,\Div^d_{\cY}}:=LG/L^+G
\]
以及 \emph{局部 Hecke 叠}
\[
\mathrm{Hk}_{G,\Div^d_{\cY}}=L^+G\backslash\Gr_{G,\Div^d_{\cY}}=L^+G\backslash LG/L^+G.
\]
这个 $\mathrm{Hk}_{G,\Div^d_{\cY}}$ 是 $\Div^d_{\cY}$ 上的小 $v$-叠。
\end{definition}

\begin{proposition}
$\Gr_{G,\Div^d_{\cY}}$ 参数化 $B^+$ 上的 $G$-主丛 $\cE$，连同其在 $B$ 上的平凡化。等价地\footnote{这一等价是「Beauville--Laszlo 粘接引理」的类比。}，它参数化 $\cY_S$ 上的 $G$-主丛 $\cE$，连同其在 $\cY_S\setminus D_S$ 上的亚纯平凡化。

$\mathrm{Hk}_{G,\Div^d_{\cY}}$ 参数化 $B^+$ 上的 $G$-主丛 $\cE_1,\cE_2$，连同张量到 $B$ 之后它们之间的一个同构。等价地，它参数化 $\cY_S$ 上的 $G$-主丛 $\cE_1,\cE_2$，连同它们限制到 $\cY_S\setminus D_S$ 之后的一个同构。
\end{proposition}

\srcpage{125}

\subsection{Schubert 簇}

对任意小 $v$-叠 $S\to\Div^d_{\cY}$，令 $\Gr_{G,S/\Div^d_{\cY}}:=\Gr_{G,\Div^d_{\cY}}\times_{\Div^d_{\cY}}S$。

例如，$S=\Spa C$ 可以是几何点，而 $S\to\Div^d_{\cY}$ 对应于 $C$ 的 $d$ 个 untilt $C^{\#}_1,\ldots,C^{\#}_d$ 的汇聚。我们允许 $C_i$ 之间有重合，但去掉重数不改变 $\Gr_{G,C/\Div^d_{\cY}}$。因此为理解纤维，不妨假设这些 untilt 互不相同。取 $\xi_1,\ldots,\xi_d\in W_{\cO_E}(\cO_C)$，使得 $\cO_{C^{\#}_i}=W_{\cO_E}(\cO_C)/(\xi_i)$。记 $\xi=\xi_1\cdots\xi_d$。

\begin{proposition}
设 $T\subset B\subset G$。则有双射
\[
X^+_*(T)^d\xrightarrow{\ \sim\ }|\mathrm{Hk}_{G,C/\Div^d_{\cY}}|,
\]
把 $(\mu_1,\ldots,\mu_d)\in X^+_*(T)^d$ 送到 $(\mu_1(\xi_1)\cdots\mu_d(\xi_d))\in LG(C)=G(B)$ 的轨道。
\end{proposition}

\begin{remark}
事实上，
\[
\mathrm{Hk}_{G,C/\Div^d_{\cY}}=\left(\prod_{i=1}^m\mathrm{Hk}_{G,C/\Div^1_{\cY}}\right)^{/S},
\]
因此为理解几何纤维，可化归到 $d=1$ 的情形。然而与经典情形不同，这些 $\mathrm{Hk}_{G,C/\Div^1_{\cY}}$ 并不全都相同，因为它们依赖于映射 $\Spa C\to\Div^1_{\cY}$，而后者由 $C^{\#}_i$ 给出。
\end{remark}

我们可以定义 $L^+G$-轨道
\[
\Gr_{G,C/\Div^d_{\cY},(\mu_1,\ldots,\mu_d)}\subset\Gr_{G,C/\Div^d_{\cY}},
\]
它们满足通常的闭包关系
\[
\Gr_{G,C/\Div^d_{\cY},\le(\mu_1,\ldots,\mu_d)}:=\ol{\Gr_{G,C/\Div^d_{\cY},(\mu_1,\ldots,\mu_d)}}=\bigcup_{(\mu'_1,\ldots,\mu'_d)\le(\mu_1,\ldots,\mu_d)}\Gr_{G,C/\Div^d_{\cY},(\mu'_1,\ldots,\mu'_d)},
\]
其中序是通常的支配序。对 $L^+G$ 取商，我们也有
\[
\mathrm{Hk}_{G,C/\Div^d_{\cY},(\mu_1,\ldots,\mu_d)}\subset\mathrm{Hk}_{G,C/\Div^d_{\cY}}
\]
及其闭包
\[
\mathrm{Hk}_{G,C/\Div^d_{\cY},\le(\mu_1,\ldots,\mu_d)}\subset\mathrm{Hk}_{G,C/\Div^d_{\cY}}.
\]

我们也可以对族定义这些，即逐纤维地应用先前的定义。在 $S=(\Div^1_{\cY})^d\to\Div^d_{\cY}$ 上，对 $\mu_1,\mu_2,\ldots,\mu_d\in X^+_*$，可以定义
\[
\Gr_{G,S/\Div^d_{\cY},(\mu_1,\ldots,\mu_d)}\subset\Gr_{G,S/\Div_{\cY}}
\]
以及
\[
\Gr_{G,S/\Div^d_{\cY},\le(\mu_1,\ldots,\mu_d)}\subset\Gr_{G,S/\Div_{\cY}}.
\]

\begin{warning}
当 untilt 重合时，必须在定义中把相应的 $\mu_i$ 相加。
\end{warning}

\begin{proposition}[{[Berk]}]
闭子叠 $\Gr_{G,S/\Div^d_{\cY},\le(\mu_1,\ldots,\mu_d)}\subset\Gr_{G,S/\Div^d_{\cY}}$ 在 $S$ 上紧合、在空间菱形中可表，且 $\dim.\mathrm{tr.g.}$ 有限。
\end{proposition}

\begin{remark}
目前还不知道显式的 pro-étale 卡片！基变换到 $(\Spa E)^{\diamond,d}=\Div^d_{\cY}$ 之后，可以给出更显式的证明 [Bence Hevesi 的硕士论文]。然而在上述一般情形（包含向特征 $p$ 的退化），已知的唯一论证是间接的，用的是某一版本的 Artin 判别法。
\end{remark}

\srcpage{126}

\begin{proposition}
在对角线之外的开 Schubert 胞腔上，$L^+G$-作用是传递的。
\end{proposition}

用 Tannaka 形式化可化归到 $\mathrm{GL}_n$ 的情形。在 $\mathrm{GL}_n$ 的情形，它归结为关于格的一个显式陈述，可以直接分析。

\begin{corollary}
$\mathrm{Hk}_{G,S/\Div^d_{\cY}}$ 的分层在对角线之外形如 $[S/\text{某个大群}]$。这个群是有限维上同调光滑群（如 $G^\diamond$）被一个无穷维「幂幺群」（例如 $\ker(L^+G\to G^\diamond)$）的扩张。

因此，在 $D_{\et}$ 的层面上，所有分层都表现得像 Artin $v$-叠（也就是说，可以忘掉幂幺部分）。事实上，$L^+G$ 在每个 $\Gr_{G,S/\Div^d_{\cY},\le(\mu_1,\ldots,\mu_d)}$ 上的作用经过一个商 $(L^+G)_{\le(\mu_1,\ldots,\mu_d)}$ 分解，后者是有限维且上同调光滑的，其核是「幂幺的」。所以
\[
D_{\et}(\mathrm{Hk}_{G,S/\Div^d_{\cY},\le\mu_\bullet},\Lambda)\cong D_{\et}((L^+G)_{\le\mu_\bullet}\backslash\Gr_{G,S/\Div^d_{\cY},\le\mu_\bullet},\Lambda).
\]
\end{corollary}

\begin{definition}
我们定义
\[
D_{\et}(\mathrm{Hk}_{G,S/\Div^d_{\cY}},\Lambda)^{bd}:=\bigcup_\mu D_{\et}(\mathrm{Hk}_{G,S/\Div^d_{\cY},\le\mu_\bullet},\Lambda)\subset D_{\et}(\mathrm{Hk}_{G,S/\Div^d_{\cY},\le\mu_\bullet},\Lambda).
\]
范畴 $D_{\et}(\mathrm{Hk}_{G,S/\Div^d_{\cY}},\Lambda)^{bd}$ 在卷积下是幺半的。其定义如下。考虑交换图
\[
\begin{tikzcd}
\mathrm{Hk}\times\mathrm{Hk} & L^+G\backslash LG\times^{L^+G}LG/L^+G \arrow[l,"\pi"'] \arrow[d,"m"] \\
& L^+G\backslash LG/L^+G=\mathrm{Hk}
\end{tikzcd}
\]
$A,B\in D_{\et}(\mathrm{Hk},\Lambda)^{bd}$ 的卷积为
\[
A\star B:=Rm_*i^*(A\boxtimes B).
\]
\end{definition}

\begin{remark}
映射 $m$ 是 ind-紧合的，因为其纤维是仿射 Grassmann 簇，所以紧合基变换适用，从而保证结合律。
\end{remark}

\subsection{$\mathrm{Hk}$ 上的反常层}

\begin{definition}
固定 $S\to\Div^d_{\cY}$。定义
\[
{}^pD^{\le0}_{\et}(\mathrm{Hk}_{G,S/\Div^d_{\cY},\le\mu},\Lambda)^{bd}\subset D(\mathrm{Hk}_{G,S/\Div^d_{\cY},\le\mu},\Lambda)^{bd}
\]
为所有 $A\in D(\mathrm{Hk}_{G,S/\Div^d_{\cY},\le\mu},\Lambda)^{bd}$ 构成的满子范畴，使得对一切几何点 $\Spa(C,C^+)\to S$，以及一切 $\mu_1,\ldots,\mu_m\in X^+_*$（其中 $m$ 是 $\Spa(C,C^+)\to S\to\Div^d_{\cY}$ 所对应的 $\Spa(C,C^+)$ 的不同 untilt 的个数），有
\[
A|_{\mathrm{Hk}_{G,\Spa(C,C^+)/\Div^d_{\cY},\le\mu}}\in D^{\le-d(\mu)}(\Lambda),
\]
其中 $d(\mu)=\sum_{i=1}^m\langle2\rho,\mu_i\rangle$。（这是「相对反常性」的直接类比。）我们定义 ${}^pD^{\ge0}$ 为 ${}^pD^{\le0}[1]$ 的右正交补。
\end{definition}

\begin{theorem}
上述定义给出一个 $t$-结构。此外，$A\in{}^pD^{\ge0}$ 当且仅当对每个几何点 $\Spa(C,C^+)\to S$，都有 $A|_{\mathrm{Hk}_{G,\Spa(C,C^+)/\Div^d_{\cY},\le\mu}}\in{}^pD^{\ge0}$。这又等价于：到所有 Schubert 胞腔 $\mathrm{Hk}_{G,\Spa(C,C^+)/\Div^d_{\cY},\le\mu}$ 的 $!$-限制落在 $D^{\ge-d(\mu)}$ 中。特别地，沿 $S'\to S$ 的拉回是 $t$-正合的。
\end{theorem}

\srcpage{127}

\begin{warning}
对概形，类似结果用的是邻近轮的反常性。然而在 $p$ 进几何中，若没有较强的假设（例如 Zariski 可构造性），这一点不成立。与之密切相关的结果 Artin 消没，在 $p$ 进几何中也不成立。
\end{warning}

\begin{example}
取 $\mathbb{A}^1_{\cO_C}$ 的 $p$ 进完备化。其一般纤维是 $\mathbb{B}_C$，特殊纤维是 $\mathbb{A}^1_k$。Artin 消没会断言：对 $\mathbb{B}_C$ 上任意反常层 $A$，都应有 $R\Gamma(\mathbb{B}_C,A)\in D^{\le1}$。但对 $A=j'_!\Lambda$ 这并不成立，其中 $j':B'_C\hookrightarrow\mathbb{B}_C$ 是半径 $1/2$ 的球的嵌入，因为 $R\Gamma(\mathbb{B}_C,A)=R\Gamma_c(B'_C,\Lambda)=\Lambda[-2]$。

令 $i:\mathbb{A}^1_k\hookrightarrow\widehat{\mathbb{A}}^1_{\cO_C}$。我们断言 $R\psi(A)=i^*Rj_*A$ 是原点处的摩天楼层 $\Lambda[-2]$。这是因为除原点外任何点的管状邻域都不与 $B'_C$ 相交。

这使得定理 23.14 的证明成为一项挑战，因为邻近轮的良好性质在经典情形中被大量使用。
\end{example}

\emph{定理 23.14 的证明概要。} 我们使用双曲局部化。首先注意，当 $G=T$ 是环面时定理是容易的。此时 $\Gr_{T,S/\Div^d_{\cY},\le(\mu_1,\ldots,\mu_d)}\to S$ 是有限的。所以该 $t$-结构就是通常的 $t$-结构。

考虑交换图
\[
\begin{tikzcd}
& \Gr_B \arrow[dl,"q"'] \arrow[dr,"p"] \\
\Gr_G && \Gr_T
\end{tikzcd}
\]
我们定义 $CT_B:=Rp_!q^*[\langle2\rho,\nu\rangle]$。若把等变性正确地设置好，可以证明它诱导出
\[
CT_B:D_{\et}(\mathrm{Hk}_G,\Lambda)^{bd}\to D_{\et}(\mathrm{Hk}_T,\Lambda)^{bd}.
\]

\begin{lemma}[关键引理]
函子 $CT_B$ 是 $t$-正合的，并且保守。
\end{lemma}

利用引理 23.17，我们从 $T$ 上的结果推出 $G$ 上想要的结果。\hfill$\square$

\emph{引理 23.17 的证明概要。} 先化归到几何点。此时 $D_{\et}(\mathrm{Hk}_G,\Lambda)^{bd}$ 按 $D_{\et}(\mathrm{Hk}_{G,(\mu_1,\ldots,\mu_d)},\Lambda)^{bd}$ 分层。它由 $D(\Lambda)$ 生成（因为 $\mathrm{Hk}_{G,(\mu_1,\ldots,\mu_d)}$ 是一个点对一个连通群的商叠）。所以只需验证 $CT_B({}^pD^{\le0})\subset{}^pD^{\le0}$ 以及 $CT({}^pD^{\ge0})\subset{}^pD^{\ge0}$，而它们分别由标准对象 $j_{\mu!}\Lambda[d_\mu]$ 与余标准对象 $Rj_{\mu*}\Lambda[d_\mu]$ 生成。

我们还没有解释这一点，但事实证明这些层在 $\Gr_{G,S}/S$ 上是 ULA 的。存在 $X=\Gr_G$ 上的 $\mathbb{G}_m$-作用，使得 $\Gr_B=X^+$、$\Gr_T=X^0$。于是可以应用双曲局部化，而它保持 ULA 性。这说明上同调是局部常值的，从而可化归到特征 $p$ 中几何点的情形。此时我们处于 Witt 向量仿射 Grassmann 簇的设定：对 $S=(\Spa\Fq)^\diamond\to\Div^d_{\cY}$，有
\[
\Gr_{G,S/\Div^d_{\cY}}=(\Gr_G^{\mathrm{Witt}})^\diamond.
\]
此外，六种运算对概形与相联系的 $v$-层是相容的。这就化归到 Witt 向量仿射 Grassmann 簇的陈述，而在该情形结果已由 Zhu 证明。（这归结为 MV 闭链维数的通常界。）保守性则由考察 MV 闭链的几何得出。\hfill$\square$
